\section{User Stories}
\label{sec:user-stories}

User stories were elicited through a series of structured discussions between the thesis authors and their supervisors, who represent the primary stakeholders identified in Sec.~\ref{sec:stakeholders}.
Each story captures a unit of observable value from the perspective of a named actor, expressed in the conventional form \textit{``as a [actor], I want [capability] so that [benefit]''}.
The stories served as the primary elicitation artefact: they were reviewed with supervisors to confirm scope and priority, subsequently decomposed into the use cases presented in Sec.~\ref{sec:use-cases}, and finally translated into the measurable functional requirements in Tab.~\ref{tab:fr}.

Tab.~\ref{tab:us} presents a representative selection of the elicited stories.
Stories are grouped by actor and labelled for traceability; the full backlog is maintained on the project documentation site.

\begingroup
\small
\begin{longtable}{@{} p{0.10\linewidth} p{0.87\linewidth} @{}}
\caption{Representative user stories elicited through stakeholder discussions, grouped by actor.}
\label{tab:us} \\
\toprule
\textbf{ID} & \textbf{Story} \\
\midrule
\endfirsthead

\multicolumn{2}{c}{\tablename~\thetable{} --- continued} \\
\toprule
\textbf{ID} & \textbf{Story} \\
\midrule
\endhead

\midrule
\multicolumn{2}{r}{\footnotesize Continued on next page} \\
\endfoot

\bottomrule
\endlastfoot

\multicolumn{2}{l}{\textit{Researcher --- Preparation}} \\
\midrule
US-R1  & As a researcher, I want a guided step-by-step setup workflow so that I can prepare an experiment without missing critical steps. \\[2pt]
US-R2  & As a researcher, I want to create and manage reading materials with optional comprehension questions so that I can build controlled stimuli sets. \\[2pt]
US-R3  & As a researcher, I want to define reusable experiment templates so that I can run multiple sessions with consistent structure. \\[2pt]
US-R4  & As a researcher, I want to calibrate the eye tracker before each participant so that gaze data is spatially accurate for that individual. \\[2pt]
US-R5  & As a researcher, I want to select the sensing mode so that I can run demonstrations or pilot tests without physical eye-tracking hardware. \\
\midrule
\multicolumn{2}{l}{\textit{Researcher --- Live Session}} \\
\midrule
US-R6  & As a researcher, I want to observe the participant's reading view on a separate screen so that I can monitor behaviour without being visible to the participant. \\[2pt]
US-R7  & As a researcher, I want to see live gaze highlighting and a rolling attention heatmap so that I can assess reading focus in real time. \\[2pt]
US-R8  & As a researcher, I want to manually apply a typographic intervention at any point in a session so that I can act on observations the automated strategy does not capture. \\[2pt]
US-R9  & As a researcher, I want to approve or reject automated intervention proposals so that I retain control over what the participant experiences. \\[2pt]
US-R10 & As a researcher, I want to include a non-adaptive control condition so that I can compare adaptive and baseline reading experiences within the same platform. \\[2pt]
US-R11 & As a researcher, I want to compare two different decision strategies across sessions so that I can evaluate which produces better reading outcomes. \\[2pt]
\midrule
\multicolumn{2}{l}{\textit{Researcher --- Post-Session}} \\
\midrule
US-R13 & As a researcher, I want to export session data in a structured format so that I can analyse gaze behaviour and intervention effects with external tools. \\[2pt]
US-R14 & As a researcher, I want to replay a recorded session so that I can review gaze paths and intervention history after the experiment concludes. \\
\midrule
\multicolumn{2}{l}{\textit{Participant}} \\
\midrule
US-P1  & As a participant, I want typographic changes to be applied without disrupting my reading so that my natural reading behaviour is preserved for analysis. \\[2pt]
US-P2  & As a participant, I want my reading position to be preserved when the presentation changes so that I do not lose my place in the text. \\[2pt]
US-P3  & As a participant, I want to answer comprehension questions after each reading item so that my understanding of the material can be assessed. \\
\midrule
\multicolumn{2}{l}{\textit{External Module Provider}} \\
\midrule
US-E1  & As an external module provider, I want to connect to the platform over a documented protocol so that I can supply decision logic without modifying the core system. \\[2pt]
US-E2  & As an external module provider, I want to receive structured context envelopes so that I can base decisions on accurate, up-to-date gaze and session state. \\[2pt]
US-E3  & As an external module provider, I want my disconnection to be handled gracefully so that an ongoing session is not interrupted by a provider failure. \\

\end{longtable}
\endgroup
